\documentclass{template}
\title{Summary - Curves}

\begin{document}
\maketitle

\section{Curvature}

Given a curve $\gamma: I \to \mathbb{R}^n$, the curvature $\kappa$ at a point is defined as:
\[
	\kappa = \left\| \frac{\dd \boldsymbol{T}}{\dd s} \right\|
\] 
where $\boldsymbol{T}$ is the unit tangent vector and $s$ is the arc length parameter.

\begin{theorem}
	If $\gamma$ is a curve defined by a smooth function $\boldsymbol{r}(t)$, then
	in 3D space, the curvature $\kappa$ can also be computed using the formula:
	\[
		\kappa = \frac{\|\boldsymbol{r}''(t) \times \boldsymbol{r}'(t)\|}{\|\boldsymbol{r}'(t)\|^3}
	\] 
	\begin{proof}
		Here is a not rigorous proof, by the meaning of the curvature and physically.
		\[
			\boldsymbol{v} = \boldsymbol{r}'(t), \quad \boldsymbol{a} = \boldsymbol{r}''(t)
		\] 
		we can decompose the acceleration vector along the tangent and normal directions:
		\[
			\boldsymbol{a} = \boldsymbol{a}_n + \boldsymbol{a}_t
			= a_n \hat{\boldsymbol{n}} + a_t \hat{\boldsymbol{\tau}}
		\] 
		hence
		\[
			a_n = \boldsymbol{a} \times \hat{\boldsymbol{\tau}}
		\] 
		since
		\[
			\hat{\boldsymbol{\tau}} = \frac{\boldsymbol{v}}{\|\boldsymbol{v}\|}
		\] 
		we have
		\[
			\rho = \frac{\|\boldsymbol{v}\|^2}{\|\boldsymbol{a} \times \hat{\boldsymbol{\tau}}\|}
			= \frac{\|\boldsymbol{v}\|^3}{\|\boldsymbol{a} \times \boldsymbol{v}\|}
		\] 
		which is the reciprocal of the curvature $\kappa$.
		
	\end{proof}
\end{theorem}

\begin{theorem}[Curvature for Plane Cartesian Parameterized Curves]
	\[
		\kappa = \frac{|x'y'' - y'x''|}{(x'^2 + y'^2)^{3/2}}
	\] 
	\begin{proof}
		\[
		\|\langle x', y', 0 \rangle \times \langle x'', y'', 0 \rangle\|
		= 
		\begin{vmatrix} x' & y' \\ x'' & y'' \end{vmatrix}
		\] 
	\end{proof}
\end{theorem}
\begin{theorem}[Curvature for Plane Cartesian $y=f(x)$]
	\[
		\kappa = \frac{|y''|}{(1 + (y')^2)^{3/2}}
	\] 
\end{theorem}
\begin{theorem}[Curvature for Plane Cartesian Implicit Curves $F(x,y)=0$]
	\[
		\kappa = \frac{|F_y^2 F_{xx} - 2 F_x F_y F_{xy} + F_x^2 F_{yy}|}{(F_x^2 + F_y^2)^{3/2}}
	\]
	
\end{theorem}


\end{document}
